\section{Residue orbits and a conditional divisor estimate}
\label{orb:section}

We now describe the lattice of possible rows after a change of coordinates
modulo a prime power.  The size of the corresponding special-linear orbit
will control the probability of a prescribed lifted matrix.  Two divisors
measure the singularity of the residue matrix; the smaller one has bounded
expectation even after all three labels have been fixed.

Throughout this section, let $B$ be a prime, let $k\geq 1$ be an
integer, and put $h=B^k$ and $R_h=\mathbb Z/h\mathbb Z$.  We interpret
$B^k$ as zero when it occurs as an entry of a matrix over $R_h$, but as the
positive integer $h$ in counting formulas.  The row span of a matrix over
$R_h$ is the $R_h$-submodule generated by its rows.

The diagonal reduction below is the prime-power-ring form of Smith's
normal form; see Smith~\cite[Articles~12--16]{Smith1861} for the classical
integer reduction and determinantal divisors. We include the elementary
proof needed here, including singular cases.

\begin{lemma}[Diagonal form and the row lattice]
\label{orb:diagonal}
Let $C\in M_3(R_h)$ have at least one unit entry.  There are invertible
matrices $P,Q\in\operatorname{GL}_3(R_h)$ and integers
$0\leq b\leq e\leq k$ such that
\[
 C=P\,\operatorname{diag}(1,B^b,B^e)\,Q.
\]
Set $D=B^b$, $E=B^e$, and $I=DE$, all as positive integers.  The subgroup
\[
 \Lambda_C=\{v\in\mathbb Z^3:
       v\bmod h\text{ belongs to the row span of }C\}
\]
is a full lattice satisfying
\begin{equation}
 [\mathbb Z^3:\Lambda_C]=I,
 \qquad E\mathbb Z^3\subseteq\Lambda_C.
 \label{orb:lattice}
\end{equation}
Moreover, $b\geq j$, for $1\leq j\leq k$, if and only if every
$2\times2$ minor of $C$ vanishes modulo $B^j$.
\end{lemma}

\begin{proof}
Move a unit entry into the first position, scale it to $1$, and clear the
rest of its row and column by invertible operations.  In the remaining
$2\times2$ block choose a nonzero entry of least $B$-adic valuation $b$,
move it to the first position of that block, and multiply by a unit to
make it $B^b$.  Every other entry of the block is divisible by $B^b$ in
$R_h$, so row and column subtraction clear its row and column.  The final
entry is still divisible by $B^b$; scaling it by a unit makes it $B^e$
with $b\leq e\leq k$.  If the block is zero, take $b=e=k$.
This proves the stated form, including all zero-pivot cases.

Invertible row and column operations preserve the ideal generated by
the $2\times2$ minors: each new minor is a linear combination of old
minors, and the inverse operations give the reverse inclusion.  For the
diagonal matrix this ideal is $B^bR_h$.  Its image modulo $B^j$ is zero
exactly when $b\geq j$, proving the last assertion and, in particular,
showing that $b$ is independent of the diagonalization.

Write $\mathscr M$ for the row span of $C$.  Left multiplication by $P$ does not
change the row span, while multiplication on the right by $Q$ is an
automorphism of $R_h^3$.  Hence
\[
 |\mathscr M|=|R_h\times DR_h\times ER_h|
      =h\frac hD\frac hE=\frac{h^3}{I}.
\]
The reduction map identifies $\mathbb Z^3/\Lambda_C$ with $R_h^3/\mathscr M$,
which proves the index formula.  Since $D$ divides $E$, the diagonal row
span contains $ER_h^3$.  Right multiplication by $Q$ preserves
$ER_h^3$, so $\mathscr M$ also contains $ER_h^3$.  Taking inverse images proves
the containment in \eqref{orb:lattice}.
\end{proof}

The preceding description permits singular matrices: either or both of
the last two diagonal entries may be zero.  The orbit estimate must
retain these cases, because zero determinants are included in the
triangle problem.

\begin{lemma}[A special-linear orbit bound]
\label{orb:orbit}
Under the hypotheses and notation of Lemma~\ref{orb:diagonal}, define
\[
 \mathcal O_C=\{GC:G\in\operatorname{SL}_3(R_h)\}.
\]
There is an absolute constant $c_0>0$ such that
\begin{equation}
 |\mathcal O_C|\geq c_0\frac{h^8}{D^3E^2}.
 \label{orb:orbit-bound}
\end{equation}
For a uniformly chosen $G\in\operatorname{SL}_3(R_h)$, the matrix $GC$
is uniform on $\mathcal O_C$.  Every matrix in this orbit has the same
row span and determinant as $C$.
\end{lemma}

\begin{proof}
Let $\mathcal G=\operatorname{GL}_3(R_h)$, and let $\mathcal H$ be the
stabilizer of $C$ under left multiplication.  Conjugation by $P$ identifies
$\mathcal H$ with the stabilizer of $\operatorname{diag}(1,D,E)$ and
preserves determinants.  An element of this latter stabilizer has first
column exactly $(1,0,0)^{\mathsf T}$.  Its second column differs from
$(0,1,0)^{\mathsf T}$ by three entries annihilated by $D$, and its third
column differs from $(0,0,1)^{\mathsf T}$ by three entries annihilated by
$E$.  The respective annihilators have $D$ and $E$ elements.  Discarding
the invertibility condition gives
\begin{equation}
 |\mathcal H|\leq D^3E^3.
 \label{orb:stabilizer}
\end{equation}

For every unit $u\in R_h^\times$ with
$u\equiv1\pmod{h/E}$, the matrix $\operatorname{diag}(1,1,u)$
stabilizes $\operatorname{diag}(1,D,E)$.  If $e<k$, there are exactly
$E$ such units.  If $e=k$, this congruence imposes no restriction, so
there are $\varphi(h)=(1-B^{-1})E$ of them.  Consequently, in all cases,
\begin{equation}
 |\det\mathcal H|\geq (1-B^{-1})E,
 \qquad
 \frac{|\det\mathcal H|}{|R_h^\times|}\geq\frac Eh.
 \label{orb:det-image}
\end{equation}
Here $\det\mathcal H$ denotes the image of the determinant homomorphism
on $\mathcal H$.

The determinant map from $\mathcal G$ onto $R_h^\times$ has kernel
$\operatorname{SL}_3(R_h)$.  Applying the orbit--stabilizer formula in
these two groups gives the exact identity
\[
 \frac{|\mathcal O_C|}{|\mathcal G C|}
 =\frac{|\det\mathcal H|}{|R_h^\times|}.
\]
Reduction modulo $B$ shows that
\[
 |\mathcal G|
 =h^9(1-B^{-1})(1-B^{-2})(1-B^{-3})
 \geq\frac{21}{64}h^9.
\]
Together with \eqref{orb:stabilizer} and \eqref{orb:det-image}, this proves
\eqref{orb:orbit-bound}, for example with $c_0=21/64$.
Every fiber of the map $G\mapsto GC$ is a coset of its stabilizer, which
proves uniformity.  The assertions about row spans and determinants
follow directly from invertibility and $\det G=1$.
\end{proof}

We next estimate the smaller divisor $D$ for digit columns.  Only the
independence and number of digit choices matter here; their prescribed
residues may be arbitrary functions of the labels.

\begin{lemma}[Conditional digit estimate]
\label{orb:conditional}
Let $r,L$ be positive integers with $r\geq2$, $r\mid L$, and $L<B$.
Let $\mathscr L$ be a label set, and for each $\lambda\in\mathscr L$,
$1\leq i\leq3$, and $0\leq u<k$, prescribe a residue
$a_{iu}(\lambda)\in\mathbb Z/r\mathbb Z$.
Fix any three labels $\lambda_1,\lambda_2,\lambda_3$, allowing
repetitions.  Independently for every $i,j,u$, choose $d_{iju}$ uniformly
from
\[
 \{a\in\{1,\ldots,L\}:a\equiv a_{iu}(\lambda_j)\pmod r\},
\]
and define $C\in M_3(R_h)$ by
\[
 C_{ij}=\sum_{u=0}^{k-1}d_{iju}B^u\pmod h.
\]
Let $b$ and $D=B^b$ be as in Lemma~\ref{orb:diagonal}, and put
$\theta=B(r/L)^4$.  Then, for every $1\leq j\leq k$,
\begin{equation}
 \Pr(b\geq j\mid\lambda_1,\lambda_2,\lambda_3)
 \leq(r/L)^{4j}.
 \label{orb:tail}
\end{equation}
If $\theta<1$, then
\begin{equation}
 \mathbb E[D\mid\lambda_1,\lambda_2,\lambda_3]
 \leq1+\sum_{j=1}^k\theta^j
 \leq\frac1{1-\theta}.
 \label{orb:moment}
\end{equation}
\end{lemma}

\begin{proof}
Every prescribed digit set has exactly $L/r$ elements.  Since the lowest
digit belongs to $\{1,\ldots,L\}$ and $L<B$, every entry of $C$ is a
unit in $R_h$.

In addition to the three labels, condition on the complete first column
of $C$ and the first entries of its other two columns.  If $b\geq j$,
Lemma~\ref{orb:diagonal} implies, for $i,v\in\{2,3\}$,
\[
 C_{iv}\equiv C_{i1}C_{11}^{-1}C_{1v}\pmod{B^j}.
\]
The right sides are fixed under this conditioning.  Each congruence
determines the first $j$ base-$B$ digits of its entry, since all chosen
digits lie between $0$ and $B-1$.  It therefore has probability either
zero or $(r/L)^j$.  The four entries use disjoint independent digit
draws, even if their labels agree.  Their joint probability is at most
$(r/L)^{4j}$.  Averaging over the additional conditioning proves
\eqref{orb:tail}, including $j=k$.

Finally, the pointwise identity
\[
 B^b=1+\sum_{j=1}^k(B^j-B^{j-1})\,\mathbf1_{\{b\geq j\}}
\]
and \eqref{orb:tail} give \eqref{orb:moment}.
\end{proof}

For the parameters $L=r^{10}$ and $B\leq200k^2L^3$ of the digit
construction, we have
\[
 \theta\leq200k^2r^{-6}.
\]
Thus, for fixed $k$ and all sufficiently large $r$,
\begin{equation}
 \mathbb E[D\mid\lambda_1,\lambda_2,\lambda_3]\leq2
 \quad\text{for every fixed triple of labels.}
 \label{orb:uniform-moment}
\end{equation}
In particular, suppose the labels are independent and uniform in a set
of size $r^d$, and let $\mathcal E$ be the event that at least two labels
coincide.  The union bound and conditioning on the labels yield
\begin{equation}
 \mathbb E[D\mathbf1_{\mathcal E}]
 \leq2\Pr(\mathcal E)\leq6r^{-d}.
 \label{orb:collision-moment}
\end{equation}
This is the estimate needed when the digit obstruction restricts a
small determinant to triples with repeated labels.
